\documentclass[11pt,a4paper]{article}

\usepackage[T1]{fontenc}
\usepackage[osf]{XCharter}
\usepackage{amsmath}
\usepackage[margin=1in]{geometry}
\usepackage{booktabs}
\usepackage{pgfplots}
\usepackage{pgfplotstable}
\usepackage{caption}
\usepackage[colorlinks,linkcolor=black,citecolor=black,urlcolor=blue]{hyperref}
\pgfplotsset{compat=1.18}

\pgfplotstableread[col sep=comma]{data/tree-scaling.csv}\treescaling
\pgfplotstableread[col sep=comma]{data/phases-16.csv}\phases
\pgfplotstableread[col sep=comma]{data/by-position-16.csv}\bypos
\pgfplotstableread[col sep=comma]{data/merge-scaling.csv}\mergescaling

\title{Repair cost under a single list edit\\[2pt]
  \large A measurement of DCG difference for \texttt{mergeSort}}
\author{Fumola}
\date{7 September 2026}

\begin{document}
\maketitle

\begin{abstract}
\noindent
Fumola has no repair algorithm. That is deliberate while the specification for
expected behaviour is still being written, and it means repair cost cannot be
measured directly. What can be measured is the difference a small edit makes to
the demanded computation graph, which bounds what any repair would have to do:
nodes both runs name and agree on are what a repair gets for free, and
everything else it must account for. This report samples that difference for
the two phases of \texttt{mergeSort} under a single list removal. The
list-into-tree phase is stable --- the number of changed nodes stays near
constant from $n=16$ to $n=1000$ while the graph grows $62\times$, so the
changed fraction falls from $0.171$ to $0.0035$. The merge network is not:
its cost is heavier, far more variable, and tracks how many merge thunks a
single symbol names.
\end{abstract}

\section{What is measured}

Two runs are compared, identical but for one removed list cell. Each run is
reset first, so its history is its own. The comparison is
\texttt{A.Diff.nodeValsDiff}, which keys nodes by pointer and ignores edge ids
and meta times because neither affects whether a cached result can be reused.
It yields four numbers:

\begin{center}
\begin{tabular}{ll}
\toprule
\textsc{onlyLeft}  & the left run named this node; the right run did not \\
\textsc{equal}     & both named it and agree on its content \\
\textsc{notEqual}  & both named it and disagree \\
\textsc{onlyRight} & the right run named it; the left did not \\
\bottomrule
\end{tabular}
\end{center}

\noindent
Write \emph{changed} for $\textsc{onlyLeft}+\textsc{onlyRight}+\textsc{notEqual}$
and \emph{diffFraction} for changed over the larger of the two runs' node
counts. A repair that had to rebuild every changed node and could reuse every
equal one would do work proportional to \emph{changed}.

This is an upper bound on the useful measurement, not the measurement itself. A
real repair may touch fewer nodes than differ, and it may touch nodes that do
not differ in order to discover that they do not. The four numbers say what
\emph{could} be reused, which is the part the naming decisions control.

\paragraph{Sampling.} Each sample draws a fresh input seed and a fresh removal
position, both from the seeded generator rather than swept, so no structure of
a sampling grid can be mistaken for structure in the answer. Removal positions
fall in $2 \ldots n$: the head is not addressable, because the edit cursor
starts one cell in (issue \#77).

\section{The list-into-tree phase scales}

\begin{center}
\pgfplotstabletypeset[
  col sep=comma,
  columns={size,n,nodes,changed,frac,fracMedian,fracP90,fracMax,lognOverN,ratio},
  columns/size/.style={column name=$n$,int detect},
  columns/n/.style={column name=samples,int detect},
  columns/nodes/.style={column name=nodes,fixed,precision=0},
  columns/changed/.style={column name=changed,fixed,precision=1},
  columns/frac/.style={column name=mean,fixed,precision=4},
  columns/fracMedian/.style={column name=median,fixed,precision=4},
  columns/fracP90/.style={column name=p90,fixed,precision=4},
  columns/fracMax/.style={column name=max,fixed,precision=4},
  columns/lognOverN/.style={column name=$\log_2 n/n$,fixed,precision=4},
  columns/ratio/.style={column name=ratio,fixed,precision=2},
  every head row/.style={before row=\toprule,after row=\midrule},
  every last row/.style={after row=\bottomrule},
]{\treescaling}
\captionof{table}{Repair cost for list-into-tree, by input size. \emph{nodes}
and \emph{changed} are means; the four \emph{diffFraction} columns summarise its
distribution; \emph{ratio} is the mean against $\log_2 n / n$.}
\end{center}

Two things stand out. The changed fraction falls by a factor of $49$ as $n$
grows $62\times$, and it stays \emph{below} $\log_2 n / n$ at every size, with
the ratio itself improving from $0.68$ to $0.35$. Figure~\ref{fig:frac} shows
this against that reference.

\begin{figure}[htbp]
\centering
\begin{tikzpicture}
\begin{loglogaxis}[
  width=0.72\textwidth, height=0.5\textwidth,
  xlabel={input size $n$}, ylabel={diffFraction},
  log basis x=2, grid=both, grid style={gray!22},
  legend pos=south west, legend cell align=left,
  legend style={draw=gray!40,font=\small},
]
\addplot[thick,mark=*,mark size=1.7pt] table[x=size,y=frac]{\treescaling};
\addlegendentry{mean, measured}
\addplot[thick,mark=triangle*,mark size=2pt,densely dashed]
  table[x=size,y=fracP90]{\treescaling};
\addlegendentry{90th percentile}
\addplot[thick,dotted,mark=none] table[x=size,y=lognOverN]{\treescaling};
\addlegendentry{$\log_2 n / n$}
\end{loglogaxis}
\end{tikzpicture}
\caption{The changed fraction of the graph, against $\log_2 n/n$. Even the 90th
percentile stays under the reference from $n=64$ up.}
\label{fig:frac}
\end{figure}

The reason is Figure~\ref{fig:changed}: the \emph{absolute} number of changed
nodes barely moves. It goes from $13.7$ at $n=16$ to $17.5$ at $n=1000$ --- a
$1.28\times$ increase across a $62\times$ increase in size. \textsc{onlyLeft}
is exactly $4.00$ at every size measured, and all of the growth is in
\textsc{notEqual}, which rises from $9.7$ to $13.5$. Divided by $\log_2 n$ the
count \emph{falls}, from $3.42$ to $1.76$, so over this range the phase is
behaving better than $O(\log n)$ per edit rather than merely as well.

\begin{figure}[htbp]
\centering
\begin{tikzpicture}
\begin{semilogxaxis}[
  width=0.72\textwidth, height=0.44\textwidth,
  xlabel={input size $n$}, ylabel={nodes},
  log basis x=2, grid=both, grid style={gray!22},
  ymin=0, ymax=20,
  legend pos=north west, legend cell align=left,
  legend style={draw=gray!40,font=\small},
]
\addplot[thick,mark=*,mark size=1.7pt] table[x=size,y=changed]{\treescaling};
\addlegendentry{changed}
\addplot[thick,mark=square*,mark size=1.6pt,densely dashed]
  table[x=size,y=notEqual]{\treescaling};
\addlegendentry{\textsc{notEqual}}
\addplot[thick,mark=o,mark size=1.6pt,dotted]
  table[x=size,y=onlyLeft]{\treescaling};
\addlegendentry{\textsc{onlyLeft}}
\end{semilogxaxis}
\end{tikzpicture}
\caption{Changed nodes per edit, on a log size axis. Near constant:
\textsc{onlyLeft} is $4.00$ throughout, and \textsc{onlyRight} is $0$ at every
sample.}
\label{fig:changed}
\end{figure}

\section{The two phases differ}

Both phases can be separated from one pair of runs, because each is built
inside its own space, and the space is the pointer's prefix. A node of the
first reads \texttt{inputTree(3-element)}; a node of the second reads
\texttt{lazyMergeSort(merge-11)(3)}. Table~\ref{tab:phases} compares them at
$n=16$ over 12 seeds and all 15 addressable positions.

\begin{center}
\pgfplotstabletypeset[
  col sep=comma,
  columns={phase,n,nodes,changed,frac,fracMedian,fracP90,fracMax},
  columns/phase/.style={string type,column name=phase},
  columns/n/.style={column name=samples,int detect},
  columns/nodes/.style={column name=nodes,fixed,precision=1},
  columns/changed/.style={column name=changed,fixed,precision=1},
  columns/frac/.style={column name=mean,fixed,precision=3},
  columns/fracMedian/.style={column name=median,fixed,precision=3},
  columns/fracP90/.style={column name=p90,fixed,precision=3},
  columns/fracMax/.style={column name=max,fixed,precision=3},
  every head row/.style={before row=\toprule,after row=\midrule},
  every last row/.style={after row=\bottomrule},
]{\phases}
\captionof{table}{The three spaces at $n=16$. The merge network is costlier on
average and much heavier in the tail.}
\label{tab:phases}
\end{center}

The averages differ by less than a factor of two, but the tails do not. The
merge network's 90th percentile is $0.590$ against the tree's $0.274$, and its
worst case $0.758$ against $0.306$.

\section{Only one of the two phases scales}

Both phases can be sampled together, at the cost of full-demand runs, which is
why this reaches only $n=100$ where the tree phase alone reaches $n=1000$.
Sample counts are small --- 40, 30, 20 and 12 --- and the merge network's
distribution is heavy-tailed, so read Table~\ref{tab:mergescaling} as a trend
and not as an estimate.

\begin{center}
\pgfplotstabletypeset[
  col sep=comma,
  columns={size,treeChanged,treeFrac,treeRatio,mergeChanged,mergeFrac,mergeP90,mergeRatio},
  columns/size/.style={column name=$n$,int detect},
  columns/treeChanged/.style={column name=changed,fixed,precision=1},
  columns/treeFrac/.style={column name=mean,fixed,precision=4},
  columns/treeRatio/.style={column name=ratio,fixed,precision=2},
  columns/mergeChanged/.style={column name=changed,fixed,precision=1},
  columns/mergeFrac/.style={column name=mean,fixed,precision=4},
  columns/mergeP90/.style={column name=p90,fixed,precision=4},
  columns/mergeRatio/.style={column name=ratio,fixed,precision=2},
  every head row/.style={
    before row={\toprule & \multicolumn{3}{c}{tree} & \multicolumn{4}{c}{merge network}\\
      \cmidrule(lr){2-4}\cmidrule(lr){5-8}},
    after row=\midrule},
  every last row/.style={after row=\bottomrule},
]{\mergescaling}
\captionof{table}{Both phases, from the same runs. \emph{ratio} is the mean
diffFraction against $\log_2 n/n$. Samples per size: 40, 30, 20, 12.}
\label{tab:mergescaling}
\end{center}

The tree phase's ratio improves steadily, $0.74 \to 0.54 \to 0.50 \to 0.41$,
and its changed count stays near $11$ at every size. The merge network's ratio
sits near $1$ for the first three sizes --- $1.02$, $0.89$, $1.09$ --- and its
changed count grows with $n$ rather than staying flat. The drop to $0.47$ at
$n=100$ comes from twelve samples and should not be read as the trend turning;
at that sample size an expensive position simply may not have been drawn.

What is solid at every size, because it does not depend on the small samples,
is the ordering: the merge network changes two to four times as many nodes as
the tree does, and its 90th percentile runs three to five times the tree's.

\section{What the merge network's names cost}

A merge thunk is currently named for the input-tree node whose stream it is
built from --- the $S$ in \texttt{merge-}$S$. That makes a symbol's influence
proportional to how many streams descend from its tree node, which is largest
at the root. Counting over 12 seeds at $n=16$, a symbol names $5.2$ merge nodes
on average and as many as $12.8$, against $\log_2 16 = 4$: the maximum is
$3.2\times$ what an evenly spread naming would give.

Figure~\ref{fig:pos} plots the consequence. Across the 15 positions, the
correlation between a position's stream-owner count and the merge network's
diffFraction is $+0.67$.

\begin{figure}[htbp]
\centering
\begin{tikzpicture}
\begin{axis}[
  width=0.86\textwidth, height=0.46\textwidth,
  xlabel={removed position}, ylabel={diffFraction},
  ybar=0pt, bar width=5.5pt,
  xtick=data, grid=major, grid style={gray!22},
  ymin=0, legend cell align=left,
  legend style={draw=gray!40,font=\small,at={(0.5,1.04)},anchor=south,legend columns=3},
]
\addplot[fill=gray!75,draw=none] table[x=pos,y=mergeFrac]{\bypos};
\addlegendentry{merge}
\addplot[fill=gray!30,draw=none] table[x=pos,y=treeFrac]{\bypos};
\addlegendentry{tree}
\addplot[sharp plot,thick,mark=*,mark size=1.4pt,color=black,
  y filter/.expression={y/40}] table[x=pos,y=streamOwner]{\bypos};
\addlegendentry{stream-owner count ($/40$)}
\end{axis}
\end{tikzpicture}
\caption{Cost by removed position at $n=16$. Positions 3 and 12 are the
expensive ones; position 3 is the symbol whose level roots every tree over
positional symbols (issue \#72).}
\label{fig:pos}
\end{figure}

Removing cell 3 rebuilds $69.5\%$ of the merge network while changing only
$27.4\%$ of the tree. Cell 3's symbol has the highest level of the first
several hundred, so it roots every tree built over positional symbols, and its
identity propagates into every stream built from it. Two effects compound:
being a heavy stream owner, and sitting high in the tree. Positions 11 and 16
own many streams but sit low, and stay cheap at $0.182$ and $0.183$ --- which
is why the correlation is $0.67$ and not higher.

This supports a renaming already proposed: name each merge thunk for the
\emph{pair of input cells whose comparison produced it}, rather than for the
tree node the stream came from. Under that scheme the root symbol should name
$O(\log n)$ nodes like any other. This report does not test that scheme, which
is not implemented; it measures the one that is, and finds the concentration
the proposal predicts.

\section{Limits}

\begin{itemize}\itemsep2pt
\item \textbf{No repair algorithm exists}, so nothing here measures repair. It
  measures what a repair could reuse. A repair that walks nodes to discover
  they are unchanged pays costs this does not capture.
\item \textbf{One edit shape.} Only removal, and only one cell at a time.
  Insertion and replacement are not sampled, and neither are edit
  \emph{sequences} --- $K$ successive edits, where reuse across edits is the
  thing that matters.
\item \textbf{The head is not addressable}, so position 1 is absent from every
  sample (issue \#77).
\item \textbf{Merge-network scaling is under-powered.} It reaches only
  $n=100$, on 40, 30, 20 and 12 samples, because it needs full-demand runs
  where the tree phase alone can skip layout. The ordering between the phases
  is solid at every size; the shape of the merge network's scaling curve is
  not established by this data.
\item \textbf{Positional symbols throughout.} Every input here is keyed
  $1 \ldots n$, which is one fixed realisation of the level function rather
  than a sample of them (issue \#72). Entropy-carrying symbols would change
  which positions are expensive, though not that some are.
\end{itemize}

\end{document}
